\subsection{導関数と増減}
平均値の定理を使えば,導関数の符号から関数の増減が分かる.まず言葉を決めておこう.区間内の任意の2点 $x_1<x_2$ についてつねに $f(x_1)<f(x_2)$ となるとき,\,$f$ はその区間で\textbf{狭義に増加する}(狭義単調増加)という.つねに $f(x_1)>f(x_2)$ となるときは\textbf{狭義に減少する}という.そのどちらかであることを\textbf{狭義単調}という.\label{用語:狭義単調}「狭義」は,等号 $f(x_1)=f(x_2)$ を許さないという意味である.

区間内の任意の2点 $x_1<x_2$ に平均値の定理を使うと,その間のある $c$ に対して
\[
 f(x_2)-f(x_1)=f'(c)(x_2-x_1)
\]
となる.\,$x_2-x_1>0$ だから,区間内で常に $f'>0$ なら $f(x_2)>f(x_1)$ であり,\,$f$ は狭義に増加する.\,$f'<0$ なら狭義に減少する.\,$f'\ge0$ の場合は $f(x_2)\ge f(x_1)$ までしか言えず,値が一定の部分もあり得る.

同じ議論から,次の結果も得られる.
\begin{theorem}{導関数が同じ関数の差}\label{定理:導関数同じ}
1つの区間のすべての内点で $F'=G'$ が成り立ち,端点を含む場合にはそこで両関数が連続なら,その区間で $F-G$ は定数である.特に,導関数が恒等的に $0$ なら,元の関数は定数である.
\end{theorem}

実際,\,$H=F-G$ と置くと $H'=0$ であり,上の式から任意の2点について $H(x_2)-H(x_1)=0$ となる.後の積分では,微分して同じ関数になる候補が,定数の違いを除いて同じであることをこの定理で確かめる.

ここで「1つの区間」という条件は外せない.たとえば
\[
 H(x)=
 \begin{cases}
  0 & (x<0)\\
  1 & (x>0)
 \end{cases}
\]
と定めれば,定義域のすべての点で $H'=0$ だが,\,$H$ は定数ではない.定義域が $x<0$ と $x>0$ の2つの区間に分かれているので,それぞれの区間で別の定数をとれるのである.
